\documentclass{article}
\usepackage[a4paper,margin=1in]{geometry}
\usepackage{graphicx}
\usepackage{enumitem}
\usepackage{amsmath}

\title{CSE411 Assignment \#3: LL(1) Top-Down Parser}
\author{Sungho Kim (20241065)}
\date{\today}

\begin{document}
\maketitle

\section{The Parsing Table for \textit{valid-grammar.json}}

\begin{table}[htbp]
\centering
\resizebox{\textwidth}{!}{%
\begin{tabular}{ l | l | l | l | l | l | l }
Nonterminal & INT & + & * & ( & ) & \$ \\
\hline
E & E → T E$'$ & & & E → T E$'$ & & \\
E$'$ & & E$'$ → + T E$'$ & & & E$'$ → $\varepsilon$ & E$'$ → $\varepsilon$ \\
T & T → F T$'$ & & & T → F T$'$ & & \\
T$'$ & & T$'$ → $\varepsilon$ & T$'$ → * F T$'$ & & T$'$ → $\varepsilon$ & T$'$ → $\varepsilon$ \\
F & F → INT & & & F → ( E ) & & 
\end{tabular}}
\end{table}

\section{The Parsing Table for \textit{invalid-grammar.json}}

\begin{table}[htbp]
\centering
\resizebox{\textwidth}{!}{%
\begin{tabular}{ l | l | l | l | l | l | l }
Nonterminal & INT & + & * & ( & ) & \$ \\
\hline
E & E → T X & & & E → T X & & \\
E$'$ & & E$'$ → + T E$'$ & & & E$'$ → $\varepsilon$ & E$'$ → $\varepsilon$ \\
T & T → INT T$'$ & & & T → ( E ) T$'$ & & \\
T$'$ & & T$'$ → $\varepsilon$ & T$'$ → * T T$'$, T$'$ → $\varepsilon$ & & T$'$ → $\varepsilon$ & T$'$ → $\varepsilon$ \\
X & & X → E$'$, X → + INT E$'$ & & & X → E$'$ & X → E$'$
\end{tabular}}
\end{table}

\section{The explanation for why \textit{invalid-grammar.json} is invalid}
The parsing table for \textit{invalid-grammar.json} contains the following conflicts:

\begin{itemize}
  \item The cell in row $\text{T}'$ and column $*$ contains both
  $\text{T}' \to *\,\text{T}\,\text{T}'$ and
  $\text{T}' \to \varepsilon$.
  This is a FIRST/FOLLOW conflict: $*$ belongs to both
  $\mathrm{FIRST}(*\,\text{T}\,\text{T}')$ and
  $\mathrm{FOLLOW}(\text{T}')$.

  \item The cell in row $\text{X}$ and column $+$ contains both
  $\text{X} \to \text{E}'$ and
  $\text{X} \to +\,\mathrm{INT}\,\text{E}'$.
  This is a FIRST/FIRST conflict: $+$ belongs to both
  $\mathrm{FIRST}(\text{E}')$ and
  $\mathrm{FIRST}(+\,\mathrm{INT}\,\text{E}')$.
\end{itemize}

These conflicts mean that an LL(1) parser cannot uniquely select a
production using only the current nonterminal and one lookahead token.
Therefore, the grammar is not LL(1) and is unsuitable for this LL(1)
parser. Here, ``invalid'' refers to its suitability for LL(1) parsing,
rather than its validity as a context-free grammar.

Although parsing table conflicts do not necessarily imply ambiguity,
this grammar is also ambiguous: some strings have more than one
distinct parse tree, as demonstrated below.

\subsection{INT + INT}

The string $\mathrm{INT} + \mathrm{INT}$ can be derived using either
production of $\text{X}$.

Using $\text{X} \to \text{E}'$:
\[
\begin{aligned}
\text{E}
  &\Rightarrow \text{T}\,\text{X} \\
  &\Rightarrow \mathrm{INT}\,\text{T}'\,\text{X} \\
  &\Rightarrow \mathrm{INT}\,\text{X} \\
  &\Rightarrow \mathrm{INT}\,\text{E}' \\
  &\Rightarrow \mathrm{INT} + \text{T}\,\text{E}' \\
  &\Rightarrow \mathrm{INT} + \mathrm{INT}\,\text{T}'\,\text{E}' \\
  &\Rightarrow \mathrm{INT} + \mathrm{INT}\,\text{E}' \\
  &\Rightarrow \mathrm{INT} + \mathrm{INT}.
\end{aligned}
\]

Using $\text{X} \to +\,\mathrm{INT}\,\text{E}'$:
\[
\begin{aligned}
\text{E}
  &\Rightarrow \text{T}\,\text{X} \\
  &\Rightarrow \mathrm{INT}\,\text{T}'\,\text{X} \\
  &\Rightarrow \mathrm{INT}\,\text{X} \\
  &\Rightarrow \mathrm{INT} + \mathrm{INT}\,\text{E}' \\
  &\Rightarrow \mathrm{INT} + \mathrm{INT}.
\end{aligned}
\]

These are two distinct leftmost derivations of the same string, so
they correspond to different parse trees. Therefore, the grammar
is ambiguous.

\subsection{INT * INT * INT}

The string $\mathrm{INT} * \mathrm{INT} * \mathrm{INT}$ can also be
derived in two different ways. In the production
$\text{T}' \to *\,\text{T}\,\text{T}'$, the nested $\text{T}$ can itself
generate further multiplication through
$\text{T} \to \mathrm{INT}\,\text{T}'$.
Consequently, the final multiplication can be generated either by
the $\text{T}'$ inside this nested $\text{T}$ or by the $\text{T}'$
following it.

Using the $\text{T}'$ inside the nested $\text{T}$ to generate the
final multiplication:
\[
\begin{aligned}
\text{E}
  &\Rightarrow \text{T}\,\text{X} \\
  &\Rightarrow \mathrm{INT}\,\text{T}'\,\text{X} \\
  &\Rightarrow \mathrm{INT} * \text{T}\,\text{T}'\,\text{X} \\
  &\Rightarrow \mathrm{INT} * \mathrm{INT}\,\text{T}'\,\text{T}'\,\text{X} \\
  &\Rightarrow \mathrm{INT} * \mathrm{INT} * \text{T}\,\text{T}'\,\text{T}'\,\text{X} \\
  &\Rightarrow \mathrm{INT} * \mathrm{INT} * \mathrm{INT}\,\text{T}'\,\text{T}'\,\text{T}'\,\text{X} \\
  &\Rightarrow \mathrm{INT} * \mathrm{INT} * \mathrm{INT}\,\text{T}'\,\text{T}'\,\text{X} \\
  &\Rightarrow \mathrm{INT} * \mathrm{INT} * \mathrm{INT}\,\text{T}'\,\text{X} \\
  &\Rightarrow \mathrm{INT} * \mathrm{INT} * \mathrm{INT}\,\text{X} \\
  &\Rightarrow \mathrm{INT} * \mathrm{INT} * \mathrm{INT}\,\text{E}' \\
  &\Rightarrow \mathrm{INT} * \mathrm{INT} * \mathrm{INT}.
\end{aligned}
\]

Using the $\text{T}'$ following the nested $\text{T}$ to generate the
final multiplication:
\[
\begin{aligned}
\text{E}
  &\Rightarrow \text{T}\,\text{X} \\
  &\Rightarrow \mathrm{INT}\,\text{T}'\,\text{X} \\
  &\Rightarrow \mathrm{INT} * \text{T}\,\text{T}'\,\text{X} \\
  &\Rightarrow \mathrm{INT} * \mathrm{INT}\,\text{T}'\,\text{T}'\,\text{X} \\
  &\Rightarrow \mathrm{INT} * \mathrm{INT}\,\text{T}'\,\text{X} \\
  &\Rightarrow \mathrm{INT} * \mathrm{INT} * \text{T}\,\text{T}'\,\text{X} \\
  &\Rightarrow \mathrm{INT} * \mathrm{INT} * \mathrm{INT}\,\text{T}'\,\text{T}'\,\text{X} \\
  &\Rightarrow \mathrm{INT} * \mathrm{INT} * \mathrm{INT}\,\text{T}'\,\text{X} \\
  &\Rightarrow \mathrm{INT} * \mathrm{INT} * \mathrm{INT}\,\text{X} \\
  &\Rightarrow \mathrm{INT} * \mathrm{INT} * \mathrm{INT}\,\text{E}' \\
  &\Rightarrow \mathrm{INT} * \mathrm{INT} * \mathrm{INT}.
\end{aligned}
\]

These are also two distinct leftmost derivations of the same string.
They place the final multiplication in different parts of the parse
tree, providing another example of the grammar's ambiguity.
\end{document}